% Generated draft: requires mathematical and source review.
\documentclass[11pt]{article}
\usepackage[margin=25mm]{geometry}
\usepackage{fontspec}
\setmainfont{Latin Modern Roman}
\usepackage{amsmath,amssymb,mathtools,unicode-math}
\setmathfont{Latin Modern Math}
\usepackage{enumitem,needspace,xurl,hyperref}
\setlength{\parindent}{0pt}
\setlength{\parskip}{5pt}
\setlength{\emergencystretch}{3em}
\newcommand{\N}{\mathbb{N}}
\newcommand{\Z}{\mathbb{Z}}
\newcommand{\Q}{\mathbb{Q}}
\newcommand{\R}{\mathbb{R}}
\newcommand{\C}{\mathbb{C}}
\newcommand{\F}{\mathbb{F}}
\newcommand{\A}{\mathbb{A}}
\newcommand{\PP}{\mathbb{P}}
\newcommand{\E}{\mathbb{E}}
\newcommand{\eps}{\varepsilon}
\newcommand{\dd}{\,\mathrm{d}}
\begin{document}
\section*{UnsolvedMath ID: 2100704}
\textbf{Source problem code:} catalogue entry
\subsection{Definitions and mathematical statement}
For context, a regular Lagrangian torus fibration $\pi:M\to B$ has a local system $\Lambda$ of integral periods. Its classical torus-bundle Chern class is a degree-two cohomological obstruction with coefficients in that local system, conventionally lying in $H^2(B;\Lambda)$. Monodromy records the twisting of $\Lambda$ and is distinct from that Chern class; global action-angle coordinates involve more than merely naming one of these obstructions.

In the semiclassical integrable-system interpretation supplied by the research summary, spectral data come from the joint spectrum of a commuting family of quantum operators, typically depending on a small parameter $\hbar$. The record does not specify such operators, their symbols, the spectral region, or the precision of the spectral data.

\textbf{Problem (S. Vũ Ngọc).} Define the quantum Chern class and detect it. In the summary's intended reading, give a precise quantum analogue of the classical Duistermaat--Chern class and a procedure or theorem recovering that invariant from the relevant joint spectra. This requires stating the quantum objects, the invariant's mathematical target and equivalence conventions, and exactly what spectral information determines it. None is fixed in the supplied sentence, so no particular cohomology group or reconstruction theorem is asserted for the quantum class. The classical group above is contextual only. The supplied summary reports quantum-monodromy work but no verified complete definition and detection scheme under the exact name quantum Chern class.
\subsection{Short English statement}
Define a quantum analogue of the classical Chern obstruction for integrable torus fibrations and explain how the intended spectral data detect it. The source does not specify the quantum model or reconstruction statement.
\subsection{Sources}
\begin{itemize}
\item \emph{Open Problems, Questions, and Challenges in Finite-Dimensional Integrable Systems} (2018), Problem 7.4, attributed to S. Vũ Ngọc, \url{https://arxiv.org/abs/1804.03737}; source-TeX extraction, accessed 2026-07-29. The encoded name in the statement is normalized to the attribution identified in its research summary.
\item The summary names Duistermaat--Chern classes and work of Vũ Ngọc and collaborators on quantum monodromy, but supplies no particular publication or theorem defining and detecting the requested invariant.
\item Local registry \texttt{data/unsolvedmath/problems.json}, permanent ID $2100704$, source code catalogue entry, \emph{Open Problems in Integrable Systems --- Integrability and Quantisation}. The record labels this open and OPEN-TRIAGE, with NEEDS\_REVIEW source-status and rights qualifications.
\end{itemize}
\end{document}
